% Example: Advanced Quran typesetting with quran package
% Compile with XeLaTeX

\documentclass[12pt,a4paper]{article}
\usepackage{polyglossia}
\setmainlanguage[locale=mashriq]{arabic}
\setotherlanguage{english}
\newfontfamily\arabicfont[Script=Arabic]{Amiri}
\usepackage[trans=en,translt]{quran}

\begin{document}

% --- Range of ayahs within a surah ---
\quranayah[2][255][257]  % Ayat al-Kursi and following 2 ayahs

% --- Full page (Mushaf page number) ---
\quranpage[1]             % Page 1 of standard Mushaf

% --- Juz with specific ruku ---
\quranruku[1][1]          % First ruku of juz 1

% --- Get metadata ---
\surahname[1]             % "Al-Fatihah"
\surahnameenglish[1]      % "The Opening"
\ayahnumber[2][255]       % Returns ayah number
\surahayatcount[1]        % Total ayahs in surah 1 (returns 7)

% --- Bilingual: Arabic + English ---
\quranayah[1][1]
\par
\quranayahen[1][1]        % English translation

% --- Combine with custom formatting ---
\begin{quote}
  \arabicfont
  \quranayah[112][1][4]   % Full Surah Ikhlas
\end{quote}
\par
\textit{Translation:} \quransurahen[112]

\end{document}
